
This study set out to test whether a more carefully curated pre-training corpus produces a more generalizable language model at matched training scale. At approximately 300 billion training tokens, using the MMLU/HellaSwag generalization balance ratio as the primary metric, and comparing Pythia-6.9B (The Pile) against OLMo-7B (Dolma), the answer is: the prediction is directly refuted, and the structural reason is that the metric's denominator is insensitive at this training scale.

The corpus-quality hypothesis predicted OLMo-7B would achieve a higher MMLU/HellaSwag ratio than Pythia-6.9B. The matched-scale evaluation (both models within 0.5\% of 300 billion training tokens) finds the opposite: Pythia achieves ratio 0.5650 versus OLMo's 0.5383 (difference = $-0.0267$, 95\% CI [$-0.0447, -0.0069]$, Cohen's d = $-2.732)$. This directional refutation is consistent across all four benchmarks evaluated. The most informative result is not the direction of the ratio difference, which may reflect an architecture effect, but the HellaSwag convergence: both models achieve 0.4580 on commonsense reasoning at this training scale regardless of corpus quality. When the denominator of a ratio metric is a constant, the ratio cannot discriminate between the factors hypothesized to drive its numerator.

The methodological contributions of this work are: (1) demonstration that the MMLU/HellaSwag generalization balance ratio does not discriminate corpus curation quality at approximately 300 billion tokens in a cross-architecture comparison; (2) identification of HellaSwag denominator saturation as the structural explanation; and (3) characterization of the minimum requirements --- architecture-matched designs or temporal trajectory analysis --- for future studies drawing causal conclusions about corpus quality effects.

Future work should address the architecture confound directly through an architecture-matched comparison: two small-scale models (1--3 billion parameters) with identical architecture trained on matched subsets of The Pile and Dolma. Temporal trajectory analysis --- evaluating both models at approximately 143 billion and 300 billion tokens --- would test whether the ratio gap is stable (architecture-driven) or narrowing (scale-driven). Alternative generalization balance metrics with verified denominator sensitivity at intermediate training scales (e.g., GSM8K/HellaSwag, MMLU/WinoGrande) may retain discriminative power where the MMLU/HellaSwag ratio fails. Direct measurement of quality proxy scores (n-gram repetition, Flesch-Kincaid grade level, fastText language identification) on matched document samples from each corpus would verify that the assumed quality difference is empirically real in the dimensions relevant to downstream generalization. These directions together constitute a principled research agenda for resolving the corpus quality $\rightarrow$ generalization balance question that the present study was unable to answer due to its design limitations.

